\chapter[1980 --- Benacerraf et al.]{1980: Baruj Benacerraf, Jean Dausset, George D. Snell}
\label{ch:1980_Benacerraf}

\section*{Main Contribution}

\textit{Discovered the genetic basis of immune recognition (the MHC system), explaining how the immune system distinguishes self from non-self.}

\section{Official Citation}

for their discoveries concerning genetically determined structures on the cell surface that regulate immunological reactions

The 1980 Nobel Prize in Physiology or Medicine was awarded jointly to Baruj Benacerraf, Jean Dausset, and George D. Snell for their discoveries concerning genetically determined structures on the cell surface that regulate immunological reactions---the major histocompatibility complex (MHC). Their work established the biological and genetic basis of immune recognition, explained why transplanted tissues are rejected by the recipient's immune system, and provided the foundation for modern transplant medicine, immunogenetics, and our understanding of autoimmune disease.

\section{Background and Historical Context}

The problem of tissue transplantation is one of the oldest and most perplexing in medicine. For centuries, physicians had attempted to transplant tissues and organs from one person to another, and for centuries, these attempts had failed. The transplanted tissue was invariably rejected by the recipient's immune system, leading to inflammation, necrosis, and loss of the graft. The mechanism of transplant rejection was unknown, and the question of why some transplants were accepted while others were rejected remained a central mystery in immunology.

The study of transplant rejection was transformed by the work of Peter Gorer and George Snell in the 1930s and 1940s. Gorer, a British immunologist, discovered that certain blood group antigens (which he called antigen II, later shown to be identical to the H-2 antigens in mice) were also responsible for transplant rejection. Snell, an American geneticist at the Jackson Laboratory in Bar Harbor, Maine, demonstrated that transplant rejection in mice was controlled by a single genetic region---which he named the \emph{H-2} locus---that encoded a set of cell surface antigens. Snell's work established the principle that transplant compatibility was genetically determined and that the major barrier to successful transplantation was the presence of incompatible H-2 antigens on the surfaces of the donor and recipient cells.

Jean Dausset, a French hematologist working in Paris, made the human connection to Snell's mouse studies. In the 1950s, Dausset discovered that blood transfusion reactions---in which a recipient's immune system attacked transfused red blood cells---were caused by antibodies directed against antigens on the surface of white blood cells (leukocytes). Dausset identified a major of these antigens as HLA-A2, and he went on to define a series of human leukocyte antigens (HLA) that were encoded by a genetic region on chromosome 6---the human major histocompatibility complex, known as the HLA complex.

Baruj Benacerraf, an Argentine-American immunologist working at New York University, made a complementary discovery using a different experimental system. In the 1960s, Benacerraf studied the immune response of inbred guinea pigs to simple protein antigens and discovered that the ability to mount an immune response to a particular antigen was genetically controlled. He identified a set of immune response genes (Ir genes) that determined whether an animal would respond to a given antigen, and he showed that these genes were closely linked to the histocompatibility genes in the major histocompatibility complex.

\section{The Discovery/Contribution}

\subsection{George D. Snell and the H-2 Complex in Mice}

George Davis Snell was born in Bradford, Massachusetts, in 1903, and trained in genetics at Cornell University and the University of Texas. He joined the Jackson Laboratory in Bar Harbor, Maine, in 1935, and spent the rest of his career there, studying the genetics of tissue transplantation in mice.

Snell's approach was to use inbred strains of mice---strains that were genetically identical as a result of many generations of brother-sister mating---as experimental tools. By crossing different inbred strains and testing the ability of the offspring to accept skin grafts from various donor strains, Snell was able to map the genes responsible for transplant rejection to a specific region of chromosome 17, which he named the H-2 complex.

Snell's work demonstrated several key principles. First, transplant rejection was primarily controlled by a relatively small number of genes located in a single genetic region---the major histocompatibility complex. Second, the H-2 complex encoded a set of cell surface glycoproteins (later called MHC molecules) that were expressed on virtually all cells of the body and that played a central role in immune recognition. Third, the H-2 complex was significantly polymorphic---that is, there were many different alleles (variants) of the H-2 genes in the mouse population, and the degree of mismatch between the H-2 alleles of the donor and recipient determined the intensity of the transplant rejection response.

Snell's genetic studies laid the groundwork for the molecular characterization of the MHC, which was achieved in the 1970s and 1980s by many groups, including those of Snell himself, Jack Strominger, and others. The MHC was shown to encode two major classes of molecules: class I MHC molecules (H-2K, H-2D, and H-2L in mice; HLA-A, HLA-B, and HLA-C in humans), which are expressed on virtually all nucleated cells and present peptides from inside the cell to cytotoxic T lymphocytes; and class II MHC molecules (H-2A and H-2E in mice; HLA-DP, HLA-DQ, and HLA-DR in humans), which are expressed primarily on antigen-presenting cells (such as macrophages, dendritic cells, and B lymphocytes) and present peptides from extracellular antigens to helper T lymphocytes.

\subsection{Jean Dausset and the Human HLA System}

Jean-Baptiste Gabriel Dausset was born in Toulouse, France, in 1916, and trained in medicine at the University of Paris. During World War II, Dausset served as a physician in the French army in North Africa, where he gained extensive clinical experience with blood transfusion and the treatment of wounded soldiers. After the war, Dausset returned to Paris and became interested in the problem of blood transfusion reactions and the immunology of transplantation.

In 1952, Dausset made a landmark observation: he discovered that the sera of some multiply transfused patients contained antibodies that agglutinated (clumped) the white blood cells of a subset of donors. Dausset realized that these antibodies were directed against antigens on the surface of leukocytes---antigens that were distinct from the well-known blood group antigens (ABO and Rh). He named the most common of these antigens MAC (after the initials of the three patients whose sera first revealed the antigen) and subsequently designated it HLA-A2.

Dausset's discovery of the HLA system was followed by a massive international effort to define and catalog the full range of human histocompatibility antigens. The first International Histocompatibility Workshop, organized by Dausset in 1964, brought together immunologists from around the world to exchange sera and cells and to standardize the techniques for tissue typing. Over the following decades, a series of international workshops progressively defined the genes and antigens of the HLA system, revealing an notable degree of polymorphism. The HLA complex was shown to be the most polymorphic genetic system in the human genome, with thousands of alleles identified at the various HLA loci.

Dausset's clinical work demonstrated the practical importance of HLA matching in organ transplantation. He showed that the closer the HLA match between donor and recipient, the lower the risk of transplant rejection and the better the clinical outcome. This finding led to the development of tissue typing as a routine clinical procedure and to the establishment of organ sharing networks that match donors and recipients based on HLA compatibility.

\subsection{Baruj Benacerraf and Immune Response Genes}

Baruj Benacerraf was born in Caracas, Venezuela, in 1920, to a Sephardic Jewish family. He trained in medicine at the University of Paris and later moved to the United States, where he held positions at the New York University School of Medicine and later at Harvard Medical School and the Dana-Farber Cancer Institute.

Benacerraf's contribution to the understanding of the MHC came from a different direction than that of Snell and Dausset. While Snell and Dausset were studying the role of MHC molecules in transplant rejection, Benacerraf was studying the genetic control of the immune response to protein antigens. Using inbred strains of guinea pigs, Benacerraf demonstrated that certain strains were ``high responders'' to specific antigens (producing large amounts of antibody) while other strains were ``low responders'' (producing little or no antibody). He showed that this variation in immune responsiveness was controlled by a set of genes---which he named immune response (Ir) genes---that were linked to the major histocompatibility complex.

Benacerraf's discovery of Ir genes had a significant impact on immunology. It established the principle that the immune response is genetically regulated and that the MHC plays a role not only in transplant rejection but also in the normal immune response to foreign antigens. This principle was later extended by the discovery that MHC molecules function as antigen-presenting molecules---that is, they bind peptides derived from foreign antigens and present them to T lymphocytes, thereby initiating the adaptive immune response. The ability of MHC molecules to bind and present peptides is determined by the structure of the peptide-binding groove of the MHC molecule, which is encoded by the polymorphic MHC genes. Different MHC alleles encode molecules with different peptide-binding specificities, which explains why different individuals (with different MHC alleles) respond differently to the same antigen.

Benacerraf also made important contributions to the understanding of immunological tolerance---the mechanism by which the immune system distinguishes ``self'' from ``non-self'' and avoids attacking the body's own tissues. His work on tolerance and autoimmunity provided insights into the mechanisms of autoimmune diseases, such as rheumatoid arthritis, type 1 diabetes, and lupus, in which the immune system mistakenly attacks the body's own tissues.

\section{Impact on Modern Medicine}

The discoveries of Benacerraf, Dausset, and Snell have had a transformative impact on modern medicine. The most immediate practical consequence of their work was the development of tissue typing for organ transplantation. By matching the HLA types of donors and recipients, clinicians could predict and minimize the risk of transplant rejection, dramatically improving the success rates of kidney, liver, heart, and lung transplants. The establishment of organ sharing networks---such as the United Network for Organ Sharing (UNOS) in the United States---based on HLA matching has saved countless lives.

The understanding of the MHC has also been essential for the development of immunosuppressive drug therapies. Drugs such as cyclosporine, tacrolimus, and mycophenolate mofetil, which are used to prevent transplant rejection, were developed in part on the basis of our understanding of how MHC molecules activate T lymphocytes. More recently, the development of targeted immunosuppressive therapies---such as belatacept, which blocks the co-stimulatory signal required for T cell activation---has been informed by our understanding of the MHC and its role in immune recognition.

Beyond transplantation, the MHC has proved to be of central importance for understanding autoimmune disease. Many autoimmune diseases are strongly associated with specific HLA alleles. For example, HLA-B27 is strongly associated with ankylosing spondylitis, HLA-DR4 is associated with rheumatoid arthritis, and HLA-DQ2 and HLA-DQ8 are associated with celiac disease and type 1 diabetes. The discovery of these associations has enabled the development of genetic risk tests for autoimmune diseases and has provided insights into the molecular mechanisms by which MHC polymorphisms influence disease susceptibility.

The MHC is also of central importance for understanding infectious disease. Different MHC alleles confer different levels of susceptibility to various infectious diseases, including HIV, malaria, tuberculosis, and hepatitis B. This variation in susceptibility is thought to be driven by the co-evolution of pathogens and the host immune system, and it has implications for the development of vaccines and other public health interventions.

In the field of cancer immunotherapy, the MHC plays a critical role. The recognition that MHC molecules present tumor-specific peptides to T lymphocytes has led to the development of adoptive T cell therapy, checkpoint inhibitor therapy, and cancer vaccines. The effectiveness of checkpoint inhibitor drugs (such as pembrolizumab and nivolumab, which target the PD-1/PD-L1 pathway) depends in part on the expression of MHC molecules on tumor cells, and HLA typing is increasingly being used to predict which patients will respond to these therapies.

\section{Legacy}

The legacy of Benacerraf, Dausset, and Snell is woven into the fabric of modern immunology and medicine. Their discoveries transformed our understanding of the immune system from a largely phenomenological science into a molecular and genetic discipline. The concept that the immune response is regulated by genetically determined cell surface molecules---a concept that emerged from their work---is now a cornerstone of immunology and has implications for virtually every area of medicine, from transplantation to autoimmunity to cancer to infectious disease.

George Snell continued his work at the Jackson Laboratory until his retirement, and he remained an influential figure in immunogenetics until his death in 1997. Jean Dausset continued his work at the University of Paris and established the Human Polymorphism Study Centre (CEPH), which became one of the world's leading centers for human genetic research. Dausset died in 2009. Baruj Benacerraf continued his work at Harvard and the Dana-Farber Cancer Institute, serving as president of the Dana-Farber Cancer Institute from 1977 to 1992. He died in 2011.

The work of the three laureates of 1980 demonstrated that the immune system is not a simple, undifferentiated defense mechanism, but rather a sophisticated, genetically regulated system that depends on the interplay of multiple genes and cell surface molecules. Their discoveries have saved countless lives through improved transplantation, have enabled the diagnosis and treatment of autoimmune diseases, and have opened new avenues for the treatment of cancer and infectious disease.


\section{Modern Developments}

Benacerraf, Dausset, and Snell's MHC work has led to modern transplant immunology and vaccine design. Understanding of HLA typing has improved organ matching and reduced rejection. MHC-restricted T cell recognition has informed vaccine design, including universal influenza vaccines. Understanding of MHC polymorphism has revealed its role in disease susceptibility, including autoimmune diseases and drug hypersensitivity. HLA-directed therapies are being developed for autoimmune conditions.
